\newpage\thispagestyle{hl}
\hlbanner{Part II $\cdot$ Full-order physics}{4}{Dynamics: modal, transient, harmonic and random response}
\hlsection{The question}
\begin{hlquote}An explicit simulation runs for a week. Someone refines the mesh a little, the time step ends up one per cent above the limit, and in six hundred steps the displacements reach $10^{68}$ m. Nothing complained: the stability limit is sharp, not gradual.\end{hlquote}
\vspace{8pt}\begin{hlsum}\textbf{The stability limit of an explicit method is sharp, not gradual.} For a plane-stress bar with lumped mass the critical step is \cdyDtcLumped{} s; the response stays bounded up to $0.99\,\Delta t_c$ and grows by a factor \cdyGrowJustOver{} at $1.01\,\Delta t_c$. The chapter compares the five dynamic solvers, and checks modal truncation and frequency response against independent references. The code does not check stability itself, so the step is the user's responsibility.\end{hlsum}
\hlsection{Key concepts}
\begin{description}\setlength\itemsep{5pt}\setlength\parsep{0pt}\renewcommand{\descriptionlabel}[1]{\hspace{\labelsep}\sffamily\bfseries\color{hlnavy}#1}
\item[Natural frequencies and modes] $\bm K\bm\phi=\omega^2\bm M\bm\phi$, with mass-normalised shapes
\item[Newmark integration] implicit; the average-acceleration rule is stable for any step
\item[Central differences] explicit and cheap per step, stable only below $2/\omega_{max}$
\item[Critical time step] set by the highest frequency; a finer mesh shortens it
\item[Modal superposition] a few modes carry a low-frequency response; truncation error is checked
\item[FRF and PSD] frequency response function, and the random-load response through $|H|^2S_{in}$
\end{description}
\hlsection{After this chapter you can}
\begin{itemize}\setlength\itemsep{3pt}
\item choose among the five dynamic solvers for a given problem
\item compute and verify the critical time step of the explicit method
\item judge a modal truncation or a frequency-response result against an independent reference
\end{itemize}
\hlsection{Explained on the page, in short boxes}
\noindent\small Semi-discrete equations and initial conditions; Modes, orthogonality and mass normalisation; Damping ratio and Rayleigh damping; Implicit and explicit integration, stability and accuracy; Where the limit $\Delta t\le2/\omega_{max}$ comes from; Stability and accuracy; Why a few modes can be enough; Frequency response function and resonance; Power spectral density in one paragraph.
\newpage\thispagestyle{hl}
\noindent{\sffamily\small\color{hlgrey}Chapter 4 \quad Dynamics: modal, transient, harmonic and random response}\par
\hlfig{\hlimg{ch04_explicit}}{Maximum displacement of a bar integrated by central differences at 0.9 and 1.1 times the critical step.}
\hlsection{Headline result}
\begin{hlheadline}For a plane-stress bar with lumped mass the critical step is \cdyDtcLumped{} s. Up to $0.99\,\Delta t_c$ the response stays bounded; at $1.01\,\Delta t_c$ it grows by a factor \cdyGrowJustOver{} over the run.\end{hlheadline}
\hlsection{From the code}
\begin{lstlisting}[style=hlpy]
free = s.free_dofs
K = np.asarray(s.K.toarray() if hasattr(s.K, "toarray") else s.K)[np.ix_(free, free)]
Ml = np.asarray(s.M_lumped.toarray() if hasattr(s.M_lumped, "toarray") else s.M_lumped)[np.ix_(free, free)]
dtc = 2 / np.sqrt(sl.eigh(K, Ml, eigvals_only=True)[-1])   # critical step: 2 / omega_max
\end{lstlisting}
\hlsection{A stated limit}
\begin{hllimit}Explicit stability is not checked by the code: the step is the user's responsibility, and the code's bound uses the consistent mass.\end{hllimit}
